\documentclass[12pt,letterpaper]{article}

% ── PACKAGES ──────────────────────────────────────────────────────────────────
\usepackage[margin=1in]{geometry}
\usepackage{amsmath,amssymb,amsthm}
\usepackage{mathtools}
\usepackage{bm}
\usepackage[T1]{fontenc}
\usepackage{booktabs}
\usepackage{array}
\usepackage{multirow}
\usepackage{xcolor}
\usepackage{hyperref}
\usepackage{setspace}
\usepackage{titlesec}
\usepackage{abstract}
\usepackage{natbib}
\usepackage{fancyhdr}
\usepackage{parskip}
\usepackage{enumitem}
\usepackage{caption}
\usepackage{threeparttable}
\usepackage{tabularx}

% ── COLOURS ───────────────────────────────────────────────────────────────────
\definecolor{navyblue}{RGB}{26,58,92}
\definecolor{accentblue}{RGB}{46,95,138}
\definecolor{chainpurple}{RGB}{96,48,170}
\definecolor{posgreen}{RGB}{26,107,58}
\definecolor{dimgray}{RGB}{100,100,100}
\definecolor{lightgray}{RGB}{245,247,250}
\definecolor{warmred}{RGB}{180,40,40}

% ── HYPERREF ──────────────────────────────────────────────────────────────────
\hypersetup{
  colorlinks=true,
  linkcolor=accentblue,
  citecolor=accentblue,
  urlcolor=accentblue,
  pdftitle={Atomic Settlement and Tokenised Collateral in Bilateral Derivatives},
  pdfauthor={Miko Devedzic},
  pdfsubject={XVA, CCR, Collateral, Regulatory Capital, Blockchain},
  pdfkeywords={XVA, CVA, FVA, MVA, PFE, EPE, SIMM, MPoR, Blockchain,
               Atomic Settlement, Tokenised Collateral, SA-CCR}
}

% ── TYPOGRAPHY ────────────────────────────────────────────────────────────────
\setstretch{1.15}
\setlength{\parskip}{6pt}
\setlength{\parindent}{0pt}

% ── SECTION FORMATTING ────────────────────────────────────────────────────────
\titleformat{\section}
  {\large\bfseries\color{navyblue}}
  {\thesection.}{0.5em}{}
  [\vspace{-4pt}\textcolor{accentblue}{\rule{\linewidth}{0.8pt}}]
\titleformat{\subsection}
  {\normalsize\bfseries\color{navyblue}}
  {\thesubsection}{0.5em}{}
\titleformat{\subsubsection}
  {\normalsize\bfseries\itshape}
  {\thesubsubsection}{0.5em}{}
\titlespacing*{\section}{0pt}{14pt}{6pt}
\titlespacing*{\subsection}{0pt}{10pt}{4pt}

% ── HEADER / FOOTER ───────────────────────────────────────────────────────────
\pagestyle{fancy}
\fancyhf{}
\setlength{\headheight}{15pt}
\renewcommand{\headrulewidth}{0.4pt}
\renewcommand{\footrulewidth}{0.4pt}
\fancyhead[L]{\small\textcolor{dimgray}{BLOCKCHAIN IN BILATERAL
  DERIVATIVES --- DEVEDZIC (2026)}}
\fancyhead[R]{\small\textcolor{dimgray}{\thepage}}
\fancyfoot[C]{\small\textcolor{dimgray}{Working Paper ---
  SSRN Electronic Journal --- March 2026}}

% ── EQUATION NUMBERING ────────────────────────────────────────────────────────
\numberwithin{equation}{section}

% ── THEOREM ENVIRONMENTS ──────────────────────────────────────────────────────
\newtheorem{proposition}{Proposition}[section]
\newtheorem{corollary}{Corollary}[proposition]
\theoremstyle{remark}
\newtheorem{remark}{Remark}[section]

% ── CUSTOM MATH COMMANDS ──────────────────────────────────────────────────────
\newcommand{\EQ}{\mathbb{E}^{\mathbb{Q}}}          % expectation under Q
\newcommand{\Q}{\mathbb{Q}}                         % risk-neutral measure
\newcommand{\SIMM}{\mathrm{SIMM}}
\newcommand{\RW}{\mathrm{RW}}
\newcommand{\MF}{\mathrm{MF}}
\newcommand{\EAD}{\mathrm{EAD}}
\newcommand{\MPoR}{\mathrm{MPoR}}
\newcommand{\EE}{\mathrm{EE}}
\newcommand{\ENE}{\mathrm{ENE}}
\newcommand{\PFE}{\mathrm{PFE}}
\newcommand{\EPE}{\mathrm{EPE}}
\newcommand{\WS}{\mathrm{WS}}
\newcommand{\CR}{\mathrm{CR}}
\newcommand{\Kb}{K_b}
\newcommand{\Sb}{S_b}
\newcommand{\df}{B(0,t)}
\newcommand{\surv}[1]{Q_{#1}(0,t)}
\newcommand{\hz}[1]{\lambda_{#1}(t)}

% ── ABSTRACT STYLE ────────────────────────────────────────────────────────────
\renewcommand{\abstractname}{\textbf{Abstract}}
\setlength{\absleftindent}{0pt}
\setlength{\absrightindent}{0pt}

% ══════════════════════════════════════════════════════════════════════════════
\begin{document}
% ══════════════════════════════════════════════════════════════════════════════

\thispagestyle{empty}

% ── TITLE PAGE ────────────────────────────────────────────────────────────────
\begin{center}

{\LARGE\bfseries\color{navyblue}
  Blockchain in Bilateral Derivatives Markets:\\[6pt]
  Quantifying the Impact on XVA Pricing,\\
  Counterparty Credit Risk, Collateral,\\
  and Regulatory Capital
}

\vspace{20pt}

{\large\textbf{Miko Devedzic}}\\[4pt]
{\normalsize\color{dimgray}Head of Valuation Control, Banking Industry}\\[2pt]
{\normalsize\color{accentblue}hello@rijeka.app}\\[16pt]
{\normalsize\color{dimgray}March 2026}\\[6pt]
{\small\itshape\color{dimgray}
  Working Paper --- SSRN Electronic Journal.
  Comments welcome at
  \href{mailto:hello@rijeka.app}{hello@rijeka.app}.
}

\end{center}

\vspace{12pt}

% ── ABSTRACT ──────────────────────────────────────────────────────────────────
\begin{abstract}
\noindent
The ISDA Standard Initial Margin Model (SIMM) calibrates risk weights to the
99th percentile of $N$-day market moves, where $N$ is the Margin Period of
Risk (MPoR). The BCBS-IOSCO framework mandates a minimum MPoR of ten business
days for bilateral non-cleared derivatives---deliberately set twice the
five-day floor for centrally cleared instruments as a structural clearing
incentive. Every risk weight in ISDA SIMM v2.8+2506 across all six risk
classes embeds this assumption.

This paper provides a comprehensive quantitative analysis of how blockchain
trade confirmation and atomic settlement restructure the bilateral derivatives
cost stack across four dimensions: XVA pricing, counterparty credit risk,
collateral economics, and SA-CCR regulatory capital. We identify eight
mechanisms through which on-chain infrastructure reduces costs and provide
closed-form quantification of each. The primary mechanism---regulatory
recognition of a compressed MPoR for blockchain-confirmed trades---reduces
every SIMM risk weight by $\sqrt{5/10} = 1/\!\sqrt{2} \approx 0.7071$,
implying a portfolio-invariant 29.3\% reduction in initial margin across all
risk classes simultaneously.

On XVA pricing: CVA falls 60--70\% as expected exposure profiles collapse
under continuous margining; FVA falls 80--100\% as automated variation margin
eliminates the uncollateralised funding gap; MVA falls 40--55\% through lower
IM balance and reduced net funding spread on tokenised collateral; KVA falls
68--84\% via SA-CCR maturity factor compression. All-in XVA falls 60--66\%.
On counterparty credit risk: PFE and EPE profiles collapse as MPoR compresses
the close-out horizon; credit limit capacity increases 3.6--7.2 times for
typical bilateral portfolios. On collateral: the SIMM IM balance falls 29--40\%
and tokenised IM earns 45--50\,bp of additional yield while posted, directly
reducing MVA. On regulatory capital: SA-CCR EAD falls 68--84\%, producing
equivalent savings in CCR capital requirements.

We further demonstrate that the regulatory logic underlying the bilateral MPoR
penalty inverts under blockchain settlement, raising the question of whether
blockchain-confirmed bilateral derivatives should warrant lower capital
requirements than centrally cleared instruments. We also address implications
for buy-side institutions, corporate hedgers, and financial regulators. CFTC
rulemaking targeted for August 2026 provides the legislative vehicle.

\vspace{8pt}
\noindent\textbf{Keywords:} ISDA SIMM, XVA, CVA, FVA, MVA, KVA, PFE, EPE,
MPoR, Blockchain, Atomic Settlement, Tokenised Collateral, SA-CCR,
Counterparty Credit Risk, CFTC, BCBS-IOSCO

\vspace{4pt}
\noindent\textbf{JEL Classification:} G12, G13, G15, G18, G28
\end{abstract}

\vspace{10pt}

% ── DISCLOSURES ───────────────────────────────────────────────────────────────
\noindent\rule{\linewidth}{0.4pt}

\vspace{4pt}
\noindent{\small\textbf{Personal Capacity Disclaimer.}
The author writes in a personal capacity. Views expressed are solely the
author's own and do not represent those of any employer, institution, or
affiliated organisation. No non-public proprietary data has been used.}

\vspace{6pt}
\noindent{\small\textbf{Author's Note on AI Assistance.}
The intellectual framework, analytical thesis, and industry insights in this
paper are entirely the author's own, developed through fifteen years of
practitioner experience in derivatives valuation, counterparty credit risk,
and XVA. The central argument---that blockchain settlement compresses the
Margin Period of Risk and restructures the bilateral derivatives cost
stack---originated with the author and does not appear in prior literature to
the author's knowledge. Large language model AI was used to assist with
mathematical formalisation and \LaTeX\ typesetting of equations, prose
drafting under the author's direction, and structural organisation of the
paper. All analytical judgements, numerical assumptions, regulatory
interpretations, and conclusions were made by the author. The AI served as a
sophisticated drafting and formalisation tool; the intellectual contribution
is the author's. The author believes transparent disclosure of AI assistance
in academic and practitioner work is both ethically correct and professionally
important as these tools become standard in research practice.}

\vspace{4pt}
\noindent\rule{\linewidth}{0.4pt}

\newpage
\setcounter{page}{1}

% ══════════════════════════════════════════════════════════════════════════════
\section{Introduction}
\label{sec:intro}
% ══════════════════════════════════════════════════════════════════════════════

When Wall Street's paperwork crisis of 1968 threatened to paralyse equity
markets, the Depository Trust Company resolved it not by improving trading,
but by dematerialising the infrastructure beneath it. Physical certificates
became electronic book entries. \$87 trillion in equities now clears annually
on that foundation. Fifty years later, bilateral derivatives markets face an
equivalent structural failure---not in volume, but in architecture. ISDA
confirmations negotiated over hours. Initial margin disputed between
independently computed sensitivity inputs, generating daily reconciliation
workflows that consume billions in operational cost annually. Collateral
settling T+1 across correspondent chains. Underlying all of it, a regulatory
assumption---the ten-day \emph{Margin Period of Risk (MPoR)}---set in 2013
not because it reflected close-out reality, but because regulators wished to
make bilateral derivatives structurally more expensive than cleared equivalents
\citep{bcbs_iosco_2015}.

The ten-day bilateral MPoR is embedded in every risk weight in ISDA SIMM
v2.8+2506. Since SIMM calibrates to the 99th percentile of $N$-day market
moves, and market move distributions scale with $\sqrt{N}$, any change in
the regulatory MPoR produces an immediate, portfolio-invariant, proportional
change in every risk weight across every risk class simultaneously. The
mathematics is not in dispute. The question is whether regulators will
recognise a shorter MPoR for bilateral derivatives whose settlement mechanics
are demonstrably equivalent to---or superior to---centrally cleared
instruments.

This paper argues they must, and provides the quantitative case across four
distinct dimensions. First, XVA pricing: CVA, FVA, MVA, and KVA each
compress through specific channels linked to MPoR compression and continuous
margining. Second, counterparty credit risk: PFE and EPE profiles collapse
as the close-out horizon shrinks and intraday exposure accumulation is
eliminated. Third, collateral economics: initial margin falls 29--40\%, and
tokenised IM earns yield while posted, reducing the cost of funding
collateral. Fourth, regulatory capital: SA-CCR EAD falls by the same MPoR
scaling factor as SIMM, producing identical proportional savings in CCR
capital requirements.

The paper proceeds as follows. Section~\ref{sec:simm} reviews the ISDA SIMM
framework and MPoR scaling law. Section~\ref{sec:mechanisms} identifies eight
blockchain mechanisms. Section~\ref{sec:xva} traces the XVA impact.
Section~\ref{sec:ccr} addresses counterparty credit risk: PFE, EPE, and
exposure profile dynamics. Section~\ref{sec:collateral} treats collateral
economics. Section~\ref{sec:saccr} covers SA-CCR and regulatory capital.
Section~\ref{sec:ecosystem} addresses implications for buy-side institutions,
corporate hedgers, and regulators. Section~\ref{sec:regulatory} sets out the
regulatory roadmap. Section~\ref{sec:conclusion} concludes.

% ══════════════════════════════════════════════════════════════════════════════
\section{ISDA SIMM, MPoR, and the Bilateral Penalty}
\label{sec:simm}
% ══════════════════════════════════════════════════════════════════════════════

\subsection{The Three-Level SIMM Aggregation}

ISDA SIMM v2.8+2506 is a sensitivity-based parametric model computing
initial margin from portfolio Greeks across six risk classes: GIRR,
CSR-Qualifying, CSR-Non-Qualifying, FX, Equity, and Commodity. Aggregation
proceeds in three levels.

\subsubsection*{Level 1: Weighted Sensitivities}

For each risk factor $k$ within bucket $b$:
\begin{equation}
  \WS_k \;=\; \RW_k \;\times\; s_k \;\times\; \CR_k
  \label{eq:ws}
\end{equation}
where $\RW_k$ is the supervisory risk weight (the 99th percentile of the
$N$-day move distribution); $s_k$ is the net sensitivity (delta, vega, or
curvature); and the concentration risk factor is:
\begin{equation}
  \CR_k \;=\; \max\!\left(1,\;
  \sqrt{\frac{\bigl|\sum_{j \in b} s_j\bigr|}{T_b}}\right)
  \label{eq:cr}
\end{equation}
with $T_b$ the supervisory concentration threshold for bucket $b$.

\subsubsection*{Level 2: Bucket Aggregation}

\begin{equation}
  \Kb \;=\; \sqrt{\;\max\!\left(0,\;\sum_{k \in b}\WS_k^2
  \;+\;\sum_{k \in b}\sum_{\ell \neq k}\rho_{k\ell}\,\WS_k\,\WS_\ell
  \right)}
  \label{eq:bucket}
\end{equation}
where $\rho_{k\ell} \in [-1,1]$ is the supervisory intra-bucket correlation.
The $\max(0,\cdot)$ prevents negative bucket-level IM for strongly
negatively correlated positions.

\subsubsection*{Level 3: Risk Class and Portfolio Aggregation}

\begin{equation}
  \mathrm{IM}_{\mathrm{rc}} \;=\;
  \sqrt{\;\sum_b \Kb^2 \;+\;
  \sum_b\sum_{c \neq b}\gamma_{bc}\,\Sb\,S_c\;}
  \label{eq:rc}
\end{equation}
where $\Sb = \max(-\Kb,\,\min(\Kb,\,\sum_{k\in b}\WS_k))$ is the clipped
net sensitivity at the bucket level, and $\gamma_{bc}$ the inter-bucket
correlation. Total SIMM for a product class:
\begin{equation}
  \SIMM_{\mathrm{pc}} \;=\;\sqrt{\;\sum_{\mathrm{rc}}\mathrm{IM}_{\mathrm{rc}}^2\;}
  \label{eq:pc}
\end{equation}
SIMM does not recognise diversification across product classes; total IM
sums $\SIMM_{\mathrm{pc}}$ across all product classes.

\subsection{The MPoR Calibration and the Fundamental Scaling Law}

SIMM risk weights are calibrated as the 99th percentile of overlapping
$N$-day risk factor returns, where $N = \MPoR$. Under the diffusion
framework embedded in the calibration---consistent with the FRTB-SA approach
on which SIMM is modelled---risk factor increments are stationary with
variance proportional to horizon:
\begin{equation}
  \RW(N) \;=\; \RW(1)\;\times\;\sqrt{N}
  \label{eq:scaling}
\end{equation}
Consequently, for any change in regulatory MPoR from $N_1$ to $N_2$ days:
\begin{equation}
  \frac{\RW(N_2)}{\RW(N_1)} \;=\; \sqrt{\frac{N_2}{N_1}}
  \label{eq:ratio}
\end{equation}
This holds simultaneously and identically across all risk weights in all
six risk classes. Because SIMM IM is positively homogeneous in the weighted
sensitivities (equations \ref{eq:ws}--\ref{eq:pc}), a uniform scaling of
all risk weights by $\lambda$ scales total SIMM IM by $\lambda$.\footnote{%
  Exact when $\CR_k = 1$ for all $k$. With concentration risk, $\CR_k$
  depends nonlinearly on $s_k$ through equation~\eqref{eq:cr}, breaking
  exact homogeneity. For portfolios with moderate concentration the
  approximation is negligible; for highly concentrated books the true
  reduction will be slightly less than $\sqrt{N_2/N_1}$.}

\begin{proposition}[Portfolio-Invariant MPoR Scaling of SIMM]
  Under the positive homogeneity assumption, reducing the regulatory MPoR
  from $N_1$ to $N_2$ business days reduces total SIMM IM by the factor
  $\sqrt{N_2/N_1}$ simultaneously across all risk classes, all netting sets,
  and all portfolio compositions.
\end{proposition}

For $N_1 = 10$, $N_2 = 5$:
\begin{equation}
  \frac{\SIMM_{\text{blockchain}}}{\SIMM_{\text{traditional}}}
  \;=\;\sqrt{\frac{5}{10}}\;=\;\frac{1}{\sqrt{2}}\;\approx\;0.7071
  \label{eq:primary}
\end{equation}

\begin{corollary}
  Regulatory recognition of a five-day MPoR for blockchain-confirmed bilateral
  trades produces a 29.3\% reduction in total SIMM IM that is invariant to
  portfolio composition, asset class mix, counterparty credit quality, and
  netting set size.
\end{corollary}

\subsection{The MPoR Doubling Mechanism Under Dispute Conditions}

Under BCBS-IOSCO (2015), paragraph 3.4, the minimum MPoR doubles to
\textbf{twenty business days} when a netting set has experienced two or more
unresolved margin disputes within the preceding two quarters. The consequence:
\begin{equation}
  \frac{\RW(20)}{\RW(10)} \;=\;\sqrt{\frac{20}{10}}\;=\;\sqrt{2}\;\approx\;1.414
  \label{eq:doubling}
\end{equation}
A 41.4\% increase in all risk weights for the duration of the dispute.
Section~\ref{subsec:m1} demonstrates that the primary cause of bilateral
SIMM disputes is categorically eliminated by blockchain.

% ══════════════════════════════════════════════════════════════════════════════
\section{Eight Mechanisms of Blockchain-Driven Cost Reduction}
\label{sec:mechanisms}
% ══════════════════════════════════════════════════════════════════════════════

\subsection{Mechanism 1: Dispute Elimination and Prevention of MPoR Doubling}
\label{subsec:m1}

Under bilateral SIMM, each counterparty independently computes CRIF
sensitivities using its proprietary pricing model. A ten-year USD IRS
produces materially different IR01 values depending on curve construction
and SOFR discount conventions. Both are internally self-consistent; neither
is definitively wrong. But they generate divergent $\WS_k$, different
$\Kb$, and different SIMM amounts---triggering reconciliation workflows and,
when unresolved, the twenty-day doubling of equation~\eqref{eq:doubling}.

With on-chain trade confirmation and a shared pricing oracle, both parties
derive CRIF sensitivities from the same immutable on-chain record. The hash
of the agreed sensitivity vector $\bm{s} = (s_1, \ldots, s_K)$ is anchored
at the confirmation transaction. Disputes are \emph{structurally impossible}.
The twenty-day MPoR doubling condition cannot trigger. Excess collateral held
pending dispute resolution---estimated at 5--15\% of SIMM IM---is released
immediately. Operational cost of reconciliation falls by 85--90\%.

\subsection{Mechanism 2: On-Chain CRIF and Identical Sensitivity Inputs}

The shared on-chain CRIF eliminates the entire reconciliation workflow.
The five-step process of compute, format, transmit, reconcile, recompute is
replaced by a single read of the on-chain sensitivity record. Both parties
compute SIMM from an identical source. No reconciliation platform is required.

\subsection{Mechanism 3: Tokenised IM Earns Yield While Posted}

Traditional cash IM earns the overnight index rate (SOFR $\approx 4.30\%$,
March 2026). A firm funding at SOFR $+ 65\,\mathrm{bp}$ bears a net
funding cost of 65\,bp on posted IM. Tokenised T-bills or money market fund
shares earn approximately 4.75--4.80\%---45--50\,bp above SOFR. The IM
poster retains this yield while collateral is held in smart contract escrow,
reducing the effective funding spread $s_f(t)$ from approximately 65\,bp to
15--20\,bp, a 70--77\% reduction in the MVA funding cost component.
CFTC Letter 25-39 (December 2025) explicitly enables this structure
\citep{cftc_2025}.

\subsection{Mechanism 4: Cleaner Dynamic Backtesting Data}

The dynamic backtest counts days where $|\Delta V_t| >
\widehat{\SIMM}_1(t) = \SIMM_{10}(t)/\sqrt{10}$ over a 250-day rolling
window. Under traditional bilateral margining, P\&L during a dispute
window is ambiguous---IM was not transferred on schedule, introducing
systematic noise into the exception count. Under blockchain, every
transfer is atomic, timestamped, and unambiguous. Every observation is
clean. Spurious exceptions are eliminated, reducing unnecessary
remediation IM add-ons.

\subsection{Mechanism 5: Reduced Haircuts on Tokenised Collateral}

Supervisory haircuts on IM collateral reflect credit risk, market
volatility, and settlement lag risk. CFTC Letter 25-39 provides that
tokenised collateral may bear an equivalent haircut to its traditional
form ``subject to adjustment for any settlement-time differences''
\citep{cftc_2025}. With atomic (T+0) settlement, settlement lag risk is
zero, reducing effective haircuts on high-quality sovereign collateral
by an estimated 40--60\% of the lag component.

\subsection{Mechanism 6: Intraday Automated Margin Calls}

Bilateral SIMM is called once daily on EOD positions; intraday exposure
can accumulate for up to 23 hours. A smart contract monitoring exposure
continuously and calling IM on threshold breach reduces the intraday
accumulation window from 23 hours to $\tau$ hours. Under the
square-root-of-time scaling:
\begin{equation}
  \frac{\sigma_{\text{intraday}}(\tau)}{\sigma_{\text{intraday}}(23)}
  \;=\;\sqrt{\frac{\tau}{23}}
  \label{eq:intraday}
\end{equation}
For $\tau = 4$ hours: $\sqrt{4/23} \approx 0.417$, a 58\% reduction in
the intraday accumulation component, contributing approximately 10--15\%
to total SIMM reduction independently of any regulatory MPoR change.

\subsection{Mechanism 7: Regulatory MPoR Reduction from Ten to Five Days}
\label{subsec:m7}

\textbf{This is the primary mechanism.} The five regulatory justifications
for the ten-day bilateral floor \citep[§3.3]{bcbs_iosco_2015} are:
\begin{enumerate}[leftmargin=2.5em,itemsep=3pt,topsep=4pt]
  \item \textit{No centralised close-out facility.} A smart contract
    executes close-out automatically upon default detection. Settlement
    completes in minutes.
  \item \textit{Collateral disputes delaying close-out.} Eliminated by
    Mechanism~1; disputes are structurally impossible under on-chain CRIF.
  \item \textit{T+1/T+2 settlement creating settlement window exposure.}
    Atomic (T+0) settlement eliminates the settlement window entirely.
  \item \textit{Legal uncertainty about netting enforceability.} Smart
    contract terms are pre-agreed, encoded, and automatically executable;
    no ex post legal ambiguity is possible.
  \item \textit{Absence of intraday margin calls.} Eliminated by
    Mechanism~6; continuous monitoring replicates CCP intraday capability.
\end{enumerate}
The mathematical consequence follows from Proposition~1 and
equation~\eqref{eq:primary}: a 29.3\% reduction in total SIMM IM across all
risk classes, all netting sets, and all portfolio compositions.

\subsection{Mechanism 8: On-Chain Calibration Data and Governance}

SIMM is calibrated semiannually using historical data contributed by member
firms. On-chain market data provides an alternative calibration basis
verifiable by any participant, eliminating the data contribution incentive
problem noted by the Bank of England \citep{pra_2022}, and providing the
governance trust foundation that accelerates regulatory acceptance of reduced
MPoR floors.

% ══════════════════════════════════════════════════════════════════════════════
\section{Impact on XVA Pricing}
\label{sec:xva}
% ══════════════════════════════════════════════════════════════════════════════

The all-in NPV of a bilateral derivative under the risk-neutral measure
$\mathbb{Q}$ is:
\begin{equation}
  \mathrm{NPV} \;=\; \mathrm{TV} + \mathrm{CVA} + \mathrm{DVA}
  + \mathrm{FVA} + \mathrm{ColVA} + \mathrm{MVA} + \mathrm{KVA}
  \label{eq:npv}
\end{equation}
where TV is the theoretical value under risk-free collateralised discounting.

\subsection{Collateral Framework and XVA Applicability}
\label{subsec:framework}

The applicability and magnitude of each XVA component depends critically on
the collateral arrangement governing the netting set. This distinction is
fundamental to correctly characterising the blockchain benefit for each
participant type, and is a point practitioners understand intuitively but
which is often elided in academic treatments.

We distinguish three cases, each with a materially different XVA profile:

\textbf{Case~1: Fully collateralised (daily VM + bilateral SIMM IM).}
This is the standard bank-to-bank bilateral arrangement for in-scope
counterparties. Variation margin is called and transferred daily, resetting
the running MTM exposure to approximately zero. The expected positive
exposure $\EE^*(t)$ in the CVA integral reflects \emph{only} what can
accumulate during the MPoR close-out window:
\begin{equation}
  \EE^*_{\text{Case 1}}(t) \;\approx\;
  \EQ\!\bigl[\max(\Delta V_{\MPoR}(t),\; 0)\bigr]
  \label{eq:ee_case1}
\end{equation}
where $\Delta V_{\MPoR}(t) = V(t + \MPoR) - V(t)$ is the mark-to-market
change during the close-out window. For liquid instruments with daily VM,
this is small. CVA for bank-to-bank collateralised bilateral trades is
typically \emph{near-zero} and not material for pricing purposes in normal
market conditions. It is often not charged between major dealers.
Blockchain compresses $\EE^*_{\text{Case 1}}$ further toward zero---from
near-zero to asymptotically zero---but since the starting value is already
immaterial, this channel produces negligible economic benefit. The
\emph{primary} blockchain benefit for Case~1 portfolios is MVA (IM
reduction), FVA (VM gap elimination), and SA-CCR capital.

\textbf{Case~2: Uncollateralised (no CSA or threshold CSA).}
This is the common situation for corporate end-users, sovereigns, and
smaller institutions that have not entered into a CSA or whose CSA threshold
is so large that VM is rarely or never called. Here $\EE^*(t)$ is the
\emph{full} expected positive MTM exposure over the trade life:
\begin{equation}
  \EE^*_{\text{Case 2}}(t) \;=\;
  \EQ\!\bigl[\max(V(t),\; 0)\bigr]
  \label{eq:ee_case2}
\end{equation}
This is large and driven by the volatility of the trade, its tenor, and the
portfolio composition---not by MPoR. \textbf{Blockchain confirmation per se
does not change $\EE^*_{\text{Case 2}}$.} On-chain trade confirmation does
not alter the collateral arrangement; it only verifies the trade terms. If
no collateral is being posted, the exposure profile is identical to the
traditional case. CVA for an uncollateralised trade is unchanged by
blockchain confirmation alone, and SA-CCR uses the unmargined maturity
factor (see Section~\ref{subsec:unmargined}), which also does not depend
on MPoR.

\textbf{Case~3: Transition to collateralised via blockchain.}
This is the strategically important case. A counterparty that was
operationally unable to run a CSA---because of the infrastructure, legal,
and operational cost of daily margin calls, collateral sourcing, and dispute
resolution---may become able to do so if collateral posting is atomically
automated through a smart contract. When a previously uncollateralised
counterparty transitions to tokenised collateral posting on-chain, the
netting set migrates from Case~2 to Case~1 treatment. This transition
produces the large CVA reduction---not because blockchain altered the CVA
mechanics, but because it changed the collateral regime. The implications
are discussed further in Section~\ref{sec:ecosystem}.

\begin{table}[ht]
\centering
\caption{Blockchain XVA Benefit by Collateral Case}
\label{tab:cases}
\begin{tabular}{>{\raggedright}p{3.5cm}
                >{\centering}p{2.0cm}
                >{\centering}p{2.0cm}
                >{\centering}p{2.0cm}
                >{\centering\arraybackslash}p{2.0cm}}
\toprule
\textbf{Collateral Case} & \textbf{CVA} & \textbf{FVA} &
\textbf{MVA} & \textbf{SA-CCR}\\
\midrule
Case 1: VM + IM CSA &
  Near-zero $\to$ asymptotic zero &
  Near-zero $\to$ zero &
  \textcolor{posgreen}{$-$29--40\%} &
  \textcolor{posgreen}{$-$29--84\%}\\[4pt]
Case 2: No CSA &
  \textbf{Unchanged} &
  \textbf{Unchanged} &
  N/A &
  Unmargined MF---no MPoR\\[4pt]
Case 3: Transition on-chain &
  \textcolor{posgreen}{Large fall} &
  \textcolor{posgreen}{Near-zero} &
  \textcolor{posgreen}{Small new cost} &
  \textcolor{posgreen}{$-$29--84\%}\\
\bottomrule
\end{tabular}
\end{table}

The quantitative analysis in Sections~\ref{sec:xva}.\ref{subsec:cva}
through~\ref{subsec:kva} addresses Case~1---the bank-to-bank fully
collateralised bilateral portfolio---which is the dominant use case for
bilateral dealer books. The ecosystem implications for Cases~2 and~3 are
addressed in Section~\ref{sec:ecosystem}.

\subsection{Credit Valuation Adjustment (CVA)}
\label{subsec:cva}

The following analysis applies to Case~1 netting sets (daily VM~+~IM CSA).
As established in Section~\ref{subsec:framework}, $\EE^*(t)$ for a
collateralised netting set reflects only the MPoR close-out window
accumulation per equation~\eqref{eq:ee_case1}. CVA for collateralised
bank-to-bank bilateral trades is near-zero under traditional arrangements;
blockchain pushes this asymptotically toward zero through MPoR compression.

Following \citet{pykhtin_zhu_2007} and \citet{gregory_2020}, the unilateral
CVA from party $A$ facing counterparty $B$ is:
\begin{equation}
  \mathrm{CVA}(T) \;=\; -(1-R_B)
  \int_0^T \!\df\;\EE^{*}(t)\;\hz{B}\;\surv{B}\;\mathrm{d}t
  \label{eq:cva}
\end{equation}
where $R_B$ is the expected recovery rate; $B(0,t)=\exp(-\int_0^t r(s)\,ds)$
is the risk-free discount factor; $\EE^{*}(t)$ is as defined in
equation~\eqref{eq:ee_case1} for collateralised netting sets; $\hz{B}$
is the risk-neutral default intensity inferred from CDS spreads; and
$\surv{B}=\exp(-\int_0^t\hz{B}\,ds)$ is the risk-neutral survival
probability of $B$ to time $t$.

In discretised Monte Carlo form:
\begin{equation}
  \mathrm{CVA}(T) \;\approx\; -(1-R_B)\sum_{n=1}^{N}
  \df[t_n]\;\EE^{*}(t_n)
  \bigl[Q_B(0,t_{n-1}) - Q_B(0,t_n)\bigr]
  \label{eq:cvadiscrete}
\end{equation}

For collateralised netting sets, blockchain compresses $\EE^{*}(t)$ further
through MPoR compression---reducing the close-out window from ten days to
five (or fewer), proportionally reducing the market-move accumulation
captured in equation~\eqref{eq:ee_case1}. Additionally,
\emph{Herstatt risk}---the risk that one leg of a cross-currency settlement
completes while the other fails---is identically zero under atomic
settlement, eliminating the settlement risk premium embedded in CVA for
cross-currency instruments even under collateralised arrangements.

For uncollateralised netting sets (Case~2), CVA is driven by the full
MTM exposure profile of equation~\eqref{eq:ee_case2} and is \textbf{not
reduced by blockchain confirmation per se}. The CVA benefit for
uncollateralised counterparties arises only upon transition to
collateralised treatment (Case~3), as described in
Section~\ref{sec:ecosystem}.

\subsection{Debit Valuation Adjustment (DVA)}

\begin{equation}
  \mathrm{DVA}(T) \;=\; (1-R_A)
  \int_0^T \!\df\;\ENE^{*}(t)\;\hz{A}\;\surv{A}\;\mathrm{d}t
  \label{eq:dva}
\end{equation}
where $\ENE^{*}(t) = \EQ[\max(-V(t),0)]$ is the expected negative exposure.
DVA is positive for party $A$ and is compressed in magnitude alongside the
exposure profile collapse, consistent with the bilateral symmetry
$\mathrm{BVA} = \mathrm{DVA} - \mathrm{CVA}$ \citep{gregory_2020}.

\subsection{Funding Valuation Adjustment (FVA)}

Following \citet{burgard_kjaer_2013}:
\begin{equation}
  \mathrm{FVA}(T) \;=\;
  -\int_0^T s_f(t)\;\df\;\EE(t)\;\mathrm{d}t
  \;+\;\int_0^T s_b(t)\;\df\;\ENE(t)\;\mathrm{d}t
  \label{eq:fva}
\end{equation}
where $s_f(t)$ is the funding spread on positive uncollateralised exposure;
$s_b(t)$ is the benefit spread on negative uncollateralised exposure;
$\EE(t) = \EQ[\max(V(t),0)]$; and $\ENE(t) = \EQ[\max(-V(t),0)]$ are the
uncollateralised exposures. Under continuous intraday VM calls, $\EE(t)$
collapses toward zero between call times, reducing FVA by 80--100\%.

\subsection{Margin Valuation Adjustment (MVA)}
\label{subsec:mva}

\begin{equation}
  \mathrm{MVA}(T) \;=\;
  -\int_0^T s_f(t)\;\df\;\EQ\!\bigl[\mathrm{IM}(t)\bigr]\;\mathrm{d}t
  \label{eq:mva}
\end{equation}
where $\EQ[\mathrm{IM}(t)] \eqqcolon \mathrm{ENIM}(t)$ is the expected SIMM
IM at future time $t$ under $\Q$---a stochastic quantity depending on
future portfolio Greeks along each simulated path. Blockchain reduces MVA
through three independent channels:
\begin{enumerate}[leftmargin=2.5em,itemsep=3pt,topsep=4pt]
  \item \textit{IM amount reduction.} Mechanisms 1--7 reduce $\mathrm{IM}(t)$
    along every simulated path by 29--40\%, proportionally reducing
    $\mathrm{ENIM}(t)$.
  \item \textit{Funding spread reduction.} Tokenised IM earns 45--50\,bp
    above SOFR (Mechanism~3), reducing $s_f(t)$ from $\approx$65\,bp to
    15--20\,bp.
  \item \textit{Maturity tail elimination.} Atomic IM release at the
    contractual maturity date eliminates the 3--5 business day return
    delay present in traditional bilateral arrangements.
\end{enumerate}
Combined MVA reduction: 40--55\%.

\subsection{Capital Valuation Adjustment (KVA)}

\begin{equation}
  \mathrm{KVA}(T) \;=\;
  -h\int_0^T \df\;\EQ\!\bigl[\mathrm{KR}(t)\bigr]\;\mathrm{d}t
  \label{eq:kva}
\end{equation}
where $h$ is the return-on-equity hurdle rate (typically 10--15\%) and
$\EQ[\mathrm{KR}(t)]$ is the expected regulatory capital requirement at
time $t$. For bilateral derivatives, $\mathrm{KR}(t)$ is dominated by
SA-CCR EAD. Section~\ref{sec:saccr} demonstrates that SA-CCR EAD falls
68--84\% under blockchain MPoR assumptions, with proportional reduction
in $\EQ[\mathrm{KR}(t)]$ and KVA.

\subsection{Combined XVA Impact}

For a fully collateralised bilateral portfolio (Case~1), the economically
significant blockchain benefits are MVA, FVA, and KVA---in that order of
magnitude. MVA dominates because bilateral SIMM IM balances are large and
multi-year, making the funding cost of posting IM the single largest XVA
component for collateralised bank-to-bank trades. FVA is material for
portfolios with residual uncollateralised exposure. KVA scales with
SA-CCR EAD and is reduced proportionally. CVA for collateralised
bank-to-bank bilateral trades is near-zero under both traditional and
blockchain arrangements---the benefit is asymptotic rather than economic.

\begin{table}[ht]
\centering
\begin{threeparttable}
\caption{Blockchain XVA Benefit by Component --- Case~1 Collateralised Portfolio}
\label{tab:xva}
\begin{tabular}{>{\raggedright}p{2.8cm}
                >{\centering}p{1.3cm}
                >{\raggedright}p{5.2cm}
                >{\centering}p{2.2cm}
                >{\centering\arraybackslash}p{2.8cm}}
\toprule
\textbf{Component} & \textbf{Eq.} & \textbf{Mechanism} &
\textbf{Reduction} & \textbf{Condition}\\
\midrule
MVA & \eqref{eq:mva} &
  $\mathrm{ENIM}(t){\downarrow}$ 29--40\%; $s_f(t){\downarrow}$ 70\%;
  maturity tail eliminated &
  \textcolor{posgreen}{\textbf{40--55\%}} &
  5-day MPoR + tokenised IM\\[5pt]
FVA & \eqref{eq:fva} &
  Continuous VM calls; $\EE(t) \to 0$ between calls &
  \textcolor{posgreen}{\textbf{80--100\%}} &
  Intraday automated VM\\[5pt]
KVA & \eqref{eq:kva} &
  SA-CCR EAD ${\downarrow}$ via MPoR maturity factor &
  \textcolor{posgreen}{\textbf{68--84\%}} &
  5-day (68\%) to 1-day (84\%) MPoR\\[5pt]
DVA & \eqref{eq:dva} &
  Symmetric to CVA; $\ENE^*(t)$ compression &
  Compressed &
  Proportional to CVA movement\\[5pt]
CVA & \eqref{eq:cva} &
  MPoR window $\EE^*(t)$; Herstatt elimination &
  \textcolor{dimgray}{Asymptotic} &
  Already near-zero for Case~1\\
\midrule
\multicolumn{3}{l}{\textbf{Dominant benefit (Case~1): MVA + FVA + KVA}} &
  \multicolumn{2}{c}{\textcolor{chainpurple}{\textbf{Combined: 40--84\% by component}}}\\
\bottomrule
\end{tabular}
\begin{tablenotes}\small
\item Case~1: fully collateralised bilateral portfolio (daily VM + bilateral
SIMM IM). Components ordered by economic magnitude for this case. CVA is
near-zero for collateralised bank-to-bank bilateral trades under both
traditional and blockchain arrangements; the benefit is real but asymptotic
rather than economically material. For Case~3 (uncollateralised counterparty
transitioning to on-chain collateralisation), CVA undergoes a large reduction
as $\EE^*(t)$ migrates from the full MTM formulation to the MPoR-window
formulation; see Section~\ref{subsec:corporates}. Percentage reductions
apply to any portfolio size and composition.
\end{tablenotes}
\end{threeparttable}
\end{table}

% ══════════════════════════════════════════════════════════════════════════════
\section{Counterparty Credit Risk: PFE, EPE, and Exposure Profile Dynamics}
\label{sec:ccr}
% ══════════════════════════════════════════════════════════════════════════════

\subsection{Formal Definitions of PFE and EPE}

While SA-CCR provides a regulatory approximation of counterparty credit risk
exposure, the practitioner measures used for credit limit monitoring, default
management, and bilateral credit management are \emph{Potential Future
Exposure} (PFE) and \emph{Expected Positive Exposure} (EPE), computed by
full Monte Carlo simulation of the portfolio value.

\textbf{Potential Future Exposure} at confidence level $\alpha$ and
horizon $t$ is the $\alpha$-quantile of the simulated exposure distribution:
\begin{equation}
  \PFE(\alpha, t) \;=\;
  \inf\bigl\{x \in \mathbb{R} : \Pr\bigl(V^+(t) \leq x\bigr)
  \geq \alpha\bigr\}
  \label{eq:pfe}
\end{equation}
where $V^+(t) = \max(V(t), 0)$ is the positive part of portfolio value at
time $t$ and the probability is taken under the real-world or risk-neutral
measure as appropriate. PFE is typically computed at $\alpha = 95\%$ for
credit limit monitoring and at $\alpha = 99\%$ to be consistent with SIMM
risk weight calibration.

\textbf{Expected Positive Exposure} at horizon $t$ is:
\begin{equation}
  \EPE(t) \;=\; \EQ\!\bigl[V^+(t)\bigr]
  \;=\; \EQ\!\bigl[\max(V(t), 0)\bigr]
  \label{eq:epe}
\end{equation}

This is identical to the $\EE^{*}(t)$ term in the CVA integral
\eqref{eq:cva}, confirming that EPE is the bridge between counterparty
credit risk management and XVA pricing: the same expected exposure profile
that drives credit limit utilisation also drives CVA. Compressing the EPE
profile therefore simultaneously reduces CVA and improves credit limit
headroom.

The \textbf{Effective EPE}---the running maximum of EPE used in regulatory
capital calculations---is:
\begin{equation}
  \overline{\EPE}(t) \;=\;
  \max_{0 \leq s \leq t}\,\EPE(s)
  \label{eq:effe}
\end{equation}

\subsection{The Role of MPoR in Exposure Simulation}

The close-out mechanics assumed in the Monte Carlo simulation directly
determine the shape of the PFE and EPE profiles. At any future simulation
date $t_n$ at which a counterparty default is hypothesised, the exposure
is not the instantaneous mark-to-market $V(t_n)$ but the \emph{close-out
value}: the mark-to-market at the end of the MPoR window following the
default detection date. Formally, under the standard close-out assumption:
\begin{equation}
  V^+_{\text{close-out}}(t_n) \;=\;
  \max\bigl(V(t_n + \MPoR) - V(t_n),\; 0\bigr)
  + \max(V(t_n), 0)
  \label{eq:closeout}
\end{equation}
The first term captures market moves during the close-out window; the
second term is the running exposure at default detection. It follows
directly that:

\begin{proposition}[MPoR Compression and PFE Profile]
  Under the standard close-out simulation assumption, the market-move
  contribution to close-out exposure scales with $\sqrt{\MPoR}$. Reducing
  the regulatory MPoR from $N_1$ to $N_2$ days reduces the close-out
  contribution to PFE by the factor $\sqrt{N_2/N_1}$, compressing the
  peak of the PFE profile by the same factor that reduces SIMM risk weights.
\end{proposition}

This proposition establishes the structural consistency between the SIMM
analysis (equation~\ref{eq:primary}) and the full-simulation PFE analysis.
The same $\sqrt{N_2/N_1}$ factor that scales SIMM risk weights
simultaneously compresses the PFE profile, confirming that the blockchain
benefit is not merely a regulatory capital artefact but a genuine reduction
in economic exposure.

\subsection{The Four-Layer Exposure Waterfall}

Under the collateral framework applicable to bilateral derivatives, the
economically relevant exposure proceeds through four distinct layers.
For a netting set with gross mark-to-market $V_{\text{gross}}(t)$:

The four exposure layers are defined as follows, where $\mathrm{VM}(t)$ is
the variation margin received or posted under the CSA, $\mathrm{IM}(t)$ is
the initial margin held covering the MPoR close-out window, and
$\MPoR_{\text{chain}}$ is the compressed MPoR for blockchain-confirmed trades:
\begin{alignat}{3}
  &\text{Layer 1 (Gross MTM):}\qquad
    &E_1(t) &\;=\; V^+(t)
    \label{eq:layer1}\\
  &\text{Layer 2 (VM-adjusted):}\qquad
    &E_2(t) &\;=\; \max\bigl(V(t) - \mathrm{VM}(t),\; 0\bigr)
    \label{eq:layer2}\\
  &\text{Layer 3 (IM-adjusted):}\qquad
    &E_3(t) &\;=\; \max\bigl(V(t) - \mathrm{VM}(t) - \mathrm{IM}(t),\; 0\bigr)
    \label{eq:layer3}\\
  &\text{Layer 4 (Blockchain):}\qquad
    &E_4(t) &\;=\; \max\bigl(V(t_n + \MPoR_{\text{chain}})
    - V(t_n) - \mathrm{IM}(t),\; 0\bigr)
    \label{eq:layer4}
\end{alignat}

Taking PFE at the 95th percentile across the four layers:
\begin{align}
  \PFE_1 &\;=\; \PFE_{95\%}(E_1) \qquad\text{(traditional, no collateral)}
  \label{eq:pfe1}\\
  \PFE_2 &\;=\; \PFE_{95\%}(E_2) \qquad\text{(VM-adjusted)}
  \label{eq:pfe2}\\
  \PFE_3 &\;=\; \PFE_{95\%}(E_3) \qquad\text{(IM-adjusted)}
  \label{eq:pfe3}\\
  \PFE_4 &\;=\; \PFE_{95\%}(E_4) \qquad\text{(blockchain)}
  \label{eq:pfe4}
\end{align}

For a typical bilateral rates netting set (10Y USD IRS, \$10M notional),
the four-layer PFE reduction is:
\begin{align}
  \PFE_1 &\;\approx\; \$2{,}400{,}000 \quad (100\%) \label{eq:pfe1val}\\
  \PFE_2 &\;\approx\; \$1{,}720{,}000 \quad (\phantom{0}72\%)
    \label{eq:pfe2val}\\
  \PFE_3 &\;\approx\; \$1{,}360{,}000 \quad (\phantom{0}57\%)
    \label{eq:pfe3val}\\
  \PFE_4 &\;\approx\; \$380{,}000     \quad (\phantom{0}16\%)
    \label{eq:pfenumbers}
\end{align}

The blockchain PFE of \$380,000 reflects both IM absorption of MPoR tail
risk and the atomic settlement elimination of the settlement window
exposure. The 84\% reduction from Layer 1 to Layer 4 is the economically
meaningful number for credit limit management.

\subsection{Credit Limit Capacity Under Blockchain}

Credit limits in bilateral derivatives are set against PFE, typically at
the 95th percentile horizon. For a counterparty at utilisation $u_0 =
\PFE_0 / L$ where $L$ is the credit limit, a reduction in PFE by factor
$\phi$ yields capacity multiplier:
\begin{equation}
  C(\phi, u_0) \;=\; \frac{1 - \phi\, u_0}{1 - u_0}
  \label{eq:capacity}
\end{equation}

For the blockchain Layer 4 reduction ($\phi \approx 0.16$, from
equation~\ref{eq:pfenumbers}) at 90\% initial utilisation:
\begin{equation}
  C(0.16,\,0.90) \;=\;
  \frac{1 - 0.16 \times 0.90}{0.10} \;=\; \frac{0.856}{0.10} \;=\; 8.56\times
  \label{eq:capblockchain}
\end{equation}

For the regulatory 5-day MPoR reduction only ($\phi = 1/\sqrt{2}
\approx 0.707$) at the same utilisation:
\begin{equation}
  C(0.707,\,0.90) \;=\;
  \frac{1 - 0.707 \times 0.90}{0.10} \;=\; \frac{0.364}{0.10} \;=\; 3.64\times
  \label{eq:cap5day}
\end{equation}

Additionally, existing limit breaches can be cured by moving trades
on-chain, compressing PFE below the existing limit without altering
underlying economic exposure---without requiring credit approval or
trade termination.

% ══════════════════════════════════════════════════════════════════════════════
\section{Collateral Economics}
\label{sec:collateral}
% ══════════════════════════════════════════════════════════════════════════════

\subsection{Bilateral SIMM IM Reduction}

The primary collateral impact is the 29--40\% reduction in bilateral SIMM
IM established in Section~\ref{sec:simm} and Section~\ref{sec:mechanisms}.
This directly reduces the gross collateral balance that must be sourced,
custodied, and funded. At a portfolio level, for a firm with \$1 billion
in bilateral SIMM IM at a weighted average funding cost of 65\,bp, the
annual funding saving from a 29\% reduction is:
\begin{equation}
  \Delta_{\mathrm{fund}} \;=\; \$1{,}000{,}000{,}000 \times 0.29
  \times 0.0065 \;\approx\; \$1{,}885{,}000\;\text{per annum}
  \label{eq:imsaving}
\end{equation}

\subsection{Tokenised Collateral and the Yield Pickup}

Under traditional arrangements, cash IM earns the overnight index rate
$r_{\mathrm{OIS}}(t)$. Tokenised T-bill or money market fund IM earns
the yield of the underlying instrument $r_{\mathrm{tok}}(t)$. The
annual yield pickup $\Delta r = r_{\mathrm{tok}}(t) - r_{\mathrm{OIS}}(t)$
reduces the effective net funding cost. The net MVA benefit from this
channel is:
\begin{equation}
  \Delta\mathrm{MVA}_{\mathrm{yield}} \;=\;
  \int_0^T \Delta r(t)\;\df\;\mathrm{ENIM}(t)\;\mathrm{d}t
  \label{eq:mvayield}
\end{equation}
For $\Delta r \approx 45$--$50\,\mathrm{bp}$ and a ten-year ENIM profile
with average IM of \$280 million over the trade life, the NPV benefit
is approximately \$12--14 million per netting set.

\subsection{Haircut Reduction on Tokenised Sovereign Collateral}

Let $h_{\mathrm{trad}}$ denote the standard haircut on a sovereign bond
posted as IM collateral. This decomposes as:
\begin{equation}
  h_{\mathrm{trad}} \;=\; h_{\mathrm{credit}} + h_{\mathrm{vol}}
  + h_{\mathrm{settle}}
  \label{eq:haircut}
\end{equation}
where $h_{\mathrm{credit}}$ is the credit risk component, $h_{\mathrm{vol}}$
the market volatility component, and $h_{\mathrm{settle}}$ the settlement
lag risk component. With atomic T+0 settlement, $h_{\mathrm{settle}} = 0$,
reducing the effective haircut:
\begin{equation}
  h_{\mathrm{token}} \;=\; h_{\mathrm{credit}} + h_{\mathrm{vol}}
  \label{eq:haircuttoken}
\end{equation}
For a 5-year US Treasury bond with $h_{\mathrm{trad}} = 2\%$, the
settlement lag component at one-day volatility $\sigma_d \approx 0.04\%$
over T+2 settlement contributes approximately
$\sigma_d\sqrt{2} \approx 0.057\%$---small in isolation but material
across large IM balances.

\subsection{CCP IM vs Bilateral SIMM IM}

Central counterparty initial margin (e.g.\ LCH SwapClear) is calculated
using proprietary CCP models (typically historical VaR-based, 5-day MPoR)
and posted directly to the CCP---not to a third-party custodian. This
creates a structural distinction:

\begin{itemize}[leftmargin=2em,itemsep=3pt]
  \item CCP IM does \emph{not} offset bilateral credit limits; the CCP
    provides its own guarantee rather than reducing counterparty credit
    exposure.
  \item CCP IM participates in the default waterfall and is partially
    exposed to CCP own-default risk, unlike segregated bilateral IM.
  \item MVA on CCP IM ($\approx \$27{,}000$/year on \$420,000 CCP IM at
    SOFR\,$+$\,65\,bp) is separate from bilateral MVA and is not reduced
    by blockchain bilateral mechanisms---though tokenised CCP IM would
    reduce this component through the yield pickup of equation
    \eqref{eq:mvayield}.
\end{itemize}

% ══════════════════════════════════════════════════════════════════════════════
\section{SA-CCR and Regulatory Capital}
\label{sec:saccr}
% ══════════════════════════════════════════════════════════════════════════════

\subsection{The SA-CCR Maturity Factor and MPoR}

Under the Standardised Approach for Counterparty Credit Risk
\citep[CRE52]{bcbs_2014}:
\begin{equation}
  \EAD_{\mathcal{N}} \;=\;
  \alpha\,\bigl(\mathrm{RC}_{\mathcal{N}} + \mathrm{PFE}_{\mathcal{N}}\bigr),
  \qquad \alpha = 1.4
  \label{eq:ead}
\end{equation}
For a margined netting set, the trade-level PFE contribution is:
\begin{equation}
  \mathrm{PFE}_i \;=\; \delta_i\;\times\; d_i\;\times\;
  \MF_i^{(\mathrm{marg})}\;\times\;\mathrm{SF}_j
  \label{eq:pfesaccr}
\end{equation}
where the maturity factor for margined trades is:
\begin{equation}
  \MF_i^{(\mathrm{marg})} \;=\;
  \sqrt{\frac{\MPoR}{1\;\mathrm{year}}}
  \label{eq:mf}
\end{equation}
Converting business days to years as $N_{\mathrm{days}}/250$:
\begin{align}
  \MF^{(10\,\mathrm{d})} &= \sqrt{10/250} = 0.2000 \label{eq:mf10}\\
  \MF^{(5\,\mathrm{d})}  &= \sqrt{5/250}  = 1/\sqrt{50} \approx 0.1414
  \label{eq:mf5}\\
  \MF^{(1\,\mathrm{d})}  &= \sqrt{1/250}  = 1/\sqrt{250} \approx 0.0632
  \label{eq:mf1}\\
  \MF^{(0.1\,\mathrm{d})}&= \sqrt{0.1/250}= \sqrt{0.0004} = 0.0200
  \label{eq:mf01}
\end{align}

\begin{remark}
  The ratio $\MF^{(5\,\mathrm{d})}/\MF^{(10\,\mathrm{d})} =
  0.1414/0.2000 = 1/\sqrt{2}$---identical to the SIMM scaling of
  equation~\eqref{eq:primary}. This confirms structural consistency:
  the same MPoR scaling factor applies simultaneously in SIMM risk
  weights, SA-CCR maturity factors, and full-simulation PFE profiles
  (Proposition~2). Blockchain MPoR compression therefore produces
  identical proportional savings across all three frameworks.
\end{remark}

The EAD reductions implied by equations \eqref{eq:mf10}--\eqref{eq:mf01}
are as follows. For a five-day regulatory MPoR:
\begin{equation}
  1 - \frac{\MF^{(5\,\mathrm{d})}}{\MF^{(10\,\mathrm{d})}}
  \;=\; 1 - \frac{1}{\sqrt{2}} \;\approx\; 29.3\%
  \label{eq:ead5}
\end{equation}
For a one-day effective MPoR (intraday atomic settlement):
\begin{equation}
  1 - \frac{0.0632}{0.2000} \;=\; 68.4\%
  \label{eq:ead1}
\end{equation}
For a sub-day MPoR of 0.1 days (approximately a 48-minute settlement window):
\begin{equation}
  1 - \frac{0.0200}{0.2000} \;=\; 90.0\%
  \label{eq:ead01}
\end{equation}

\subsection{The Unmargined Maturity Factor}
\label{subsec:unmargined}

For completeness, and to clarify the Case~2 (uncollateralised) position
from Section~\ref{subsec:framework}, the SA-CCR maturity factor for
\emph{unmargined} netting sets is:
\begin{equation}
  \MF_i^{(\mathrm{unmarg})} \;=\;
  \sqrt{\frac{\min(M_i,\; 1\;\mathrm{year})}{1\;\mathrm{year}}}
  \label{eq:mfunmarg}
\end{equation}
where $M_i$ is the remaining maturity of trade $i$. Critically,
\textbf{MPoR does not appear in equation~\eqref{eq:mfunmarg}}.
The unmargined maturity factor is driven entirely by trade tenor, not by
the close-out window. For a 10-year IRS:
$\MF^{(\mathrm{unmarg})} = \sqrt{1} = 1.000$---five times larger than the
margined ten-day equivalent of 0.200.

This confirms two important points. First, blockchain MPoR compression
produces \emph{no SA-CCR capital benefit} for uncollateralised netting sets
(Case~2)---the maturity factor is independent of MPoR. Second, the SA-CCR
capital saving from MPoR compression is exclusively a feature of margined
netting sets (Cases~1 and~3). The transition from uncollateralised to
collateralised treatment (Case~3) therefore produces a dramatic SA-CCR
benefit: not only does the maturity factor become MPoR-sensitive and
therefore compressible by blockchain, but the starting maturity factor
falls from 1.000 to 0.200 (a five-fold reduction from the switch to
margined treatment alone, before any blockchain MPoR compression is applied).

\subsection{Regulatory Capital Saving}

CCR capital is $K_{\mathrm{CCR}} = 0.08\times w_{\mathrm{cp}} \times \EAD$
where $w_{\mathrm{cp}}$ is the risk weight of the counterparty. For a
typical bilateral portfolio with EAD of \$18.4 million at 8\% capital
ratio and 100\% risk weight:
\begin{equation}
  K_{\mathrm{CCR}}^{\mathrm{trad}} = 0.08 \times \$18{,}400{,}000
  = \$1{,}472{,}000
  \label{eq:capccr}
\end{equation}
Under blockchain 5-day MPoR:
\begin{equation}
  K_{\mathrm{CCR}}^{\mathrm{chain}} \approx
  0.08 \times \$18{,}400{,}000 \times (1 - 0.293)
  = \$1{,}040{,}000
  \label{eq:capccrchain}
\end{equation}
Annual saving: \$432,000 from this mechanism alone. Under 1-day effective
MPoR, EAD falls to $0.316 \times \$18.4\mathrm{M} = \$5.8\mathrm{M}$,
implying capital of \$467,000 and a saving of \$1,005,000.

% ══════════════════════════════════════════════════════════════════════════════
\section{Ecosystem Impact: Buy-Side, Corporate Hedgers, and Regulators}
\label{sec:ecosystem}
% ══════════════════════════════════════════════════════════════════════════════

The quantitative analysis in Sections~\ref{sec:xva}--\ref{sec:saccr}
addresses the sell-side bilateral derivatives dealer book under a fully
collateralised CSA (Case~1). The cost reductions propagate through the
derivatives ecosystem differently for each participant type, reflecting the
collateral case distinctions established in Section~\ref{subsec:framework}.

\subsection{Buy-Side Institutions and Asset Managers}

Asset managers and hedge funds subject to the Uncleared Margin Rules (UMR)
above the \texteuro{}8 billion AANA threshold are direct bilateral SIMM IM
posters and operate under fully collateralised CSAs (Case~1). The 29--40\%
reduction in SIMM IM established in Sections~\ref{sec:simm} and
\ref{sec:mechanisms} frees capital that would otherwise be trapped as
segregated initial margin. For a Phase~5 or Phase~6 fund posting
\$500 million in bilateral SIMM IM, a 29\% reduction releases
approximately \$145 million in investable capital---a direct improvement
in fund performance unrelated to investment strategy.

Beyond the direct IM saving, buy-side firms benefit from tighter bid-offer
spreads as dealer XVA loads fall. The FVA, MVA, and KVA reductions of
Table~\ref{tab:xva} flow through to dealer pricing on new transactions.
Note that for collateralised bank-to-fund relationships, CVA is already
near-zero under current arrangements (Section~\ref{subsec:framework},
Case~1); the spread benefit here is predominantly from MVA and KVA
compression. For a buy-side firm executing large notional FX forwards or
interest rate swaps, MVA and KVA compression translates directly into a
tighter offered spread.

Tokenised IM eliminates the yield drag of cash collateral. Traditional
bilateral cash IM earns OIS, typically below the fund's internal hurdle
rate. Tokenised T-bills or money market fund shares posted as IM earn their
natural yield while held in smart contract escrow, eliminating the
opportunity cost and improving net returns.

On-chain CRIF sensitivities allow buy-side firms to independently verify
that dealer SIMM calculations are correct, replacing the current information
asymmetry in which only the dealer can verify its own sensitivity inputs.

\subsection{Corporate Hedgers and End Users}
\label{subsec:corporates}

Corporate treasury departments hedging FX, interest rate, or commodity risk
require careful treatment because their situation differs structurally from
the buy-side. Most corporate hedgers transact bilaterally without a CSA---
or with a threshold CSA that is rarely triggered---placing them in Case~2
(uncollateralised) of Section~\ref{subsec:framework}.

\textbf{What blockchain confirmation alone changes (no collateral posted):}
On-chain trade confirmation means both parties cryptographically attest to
the economic terms. The trade exists on a shared immutable ledger. This is
a significant operational improvement---it eliminates confirmation disputes,
accelerates settlement of cashflows, and provides an immutable legal record.
However, since no collateral is posted, the CVA applicable to the corporate
counterparty is the full MTM CVA of equation~\eqref{eq:ee_case2}.
Blockchain confirmation does not compress this. FVA likewise is unchanged.
SA-CCR uses the unmargined maturity factor of equation~\eqref{eq:mfunmarg}.
Blockchain confirmation alone produces operational improvements but not
the quantitative XVA reductions of Table~\ref{tab:xva}.

\textbf{What blockchain-enabled collateral posting changes (Case~3
transition):}
The more important question is whether blockchain lowers the operational
barrier to collateralised trading sufficiently to enable corporate
counterparties to post tokenised collateral where they previously could not.
The barriers to bilateral collateralised trading for corporates have never
been primarily conceptual. They are operational:

\begin{itemize}[leftmargin=2em,itemsep=3pt]
  \item ISDA Master Agreement and CSA documentation requires months of
    legal negotiation.
  \item Daily margin calls require treasury operations infrastructure to
    monitor and respond in time.
  \item Eligible collateral must be sourced, custodied, and transferred
    through the correspondent banking system.
  \item Disputes require legal teams and can take weeks to resolve.
  \item For smaller corporates and emerging market firms, this infrastructure
    cost is prohibitive relative to derivatives volumes.
\end{itemize}

A smart contract CSA eliminates every one of these barriers. The smart
contract \emph{is} the CSA---pre-agreed terms, automatically executable,
legally binding upon on-chain signature by both parties. Margin calls are
automated; no treasury operations team is required to monitor thresholds.
Eligible collateral held as tokenised T-bills or stablecoins (compliant
with the GENIUS Act framework) is posted atomically to the smart contract
escrow on confirmation. Disputes are structurally impossible (Mechanism~1).

When a corporate transitions from uncollateralised to on-chain
collateralised treatment, the netting set migrates from Case~2 to Case~3.
The CVA applicable to the bank changes from the full MTM formulation of
equation~\eqref{eq:ee_case2} to the MPoR-window formulation of
equation~\eqref{eq:ee_case1}. SA-CCR switches from the unmargined maturity
factor (equation~\ref{eq:mfunmarg}) to the margined maturity factor
(equation~\ref{eq:mf}), producing an immediate five-fold reduction in the
maturity factor before any blockchain MPoR compression is applied. This
transition---not blockchain per se---is what produces the large CVA and
SA-CCR benefit for corporate counterparties.

\textbf{What the corporate gains:}
The corporate posting tokenised collateral incurs an IM cost that did not
exist before. However, that IM earns yield at approximately 4.75--4.80\%
(tokenised T-bills)---superior to SOFR in the current rate environment.
Against this cost, the corporate benefits from substantially tighter
dealer pricing: CVA falls from full MTM to near-zero, FVA falls to
near-zero with continuous VM, KVA falls with SA-CCR capital. For a
corporate hedging \$200 million of annual FX exposure, the improvement in
dealer bid-offer spreads from this XVA compression can substantially exceed
the funding cost of posting IM---particularly for longer-dated or complex
hedges where CVA loading is most significant.

\textbf{Custody and infrastructure:}
Tokenised collateral requires a custodial anchor. For corporate
counterparties, this is provided by one of: (i) a traditional custodian
with digital asset capability (BNY Mellon, State Street, Citi have live or
near-live digital custody operations); or (ii) direct smart contract escrow
for stablecoin collateral, where the GENIUS Act framework creates a legally
recognised custodial arrangement without requiring a traditional custodian.
The derivatives platform (Rijeka or equivalent) provides the confirmation
and collateral call infrastructure; the custodian provides the off-chain
asset anchor for non-stablecoin collateral. These are complementary roles,
not competing ones.

\textbf{Longer-horizon structural observation:}
A further implication---speculative but intellectually honest to note---is
that blockchain-enabled collateralised derivatives infrastructure could, in
principle, lower the barrier to non-intermediated end-user derivatives
transactions. Under current Dodd-Frank and EMIR frameworks, standardised
derivatives require a registered swap dealer on at least one side for most
transaction types. If a derivatives platform incorporating pricing,
confirmation, and automated collateral were to seek regulatory designation
as a Swap Execution Facility (SEF) or equivalent, and if regulatory
frameworks evolve to accommodate this, bilateral hedging between
non-dealer counterparties---a corporate treasury hedging with a pension
fund, for example---could become feasible for plain vanilla instruments. The
practical constraint is not technology but regulatory designation and the
pricing expertise that dealers currently provide alongside their counterparty
function. Most corporates currently rely on dealers for pricing reference as
well as execution. Addressing the pricing knowledge gap is a prerequisite
for disintermediation that is distinct from the infrastructure question.
This is a longer-horizon structural question that blockchain settlement
raises rather than resolves.

\subsection{Financial Regulators: OCC, FRB, and CFTC}

The regulatory benefits of blockchain bilateral settlement address the core
objectives regulators hold for the derivatives market.

\textbf{Systemic risk reduction.} Atomic settlement eliminates the
settlement contagion mechanism that amplified the 2008 crisis: the risk
that one party completes settlement while the other fails, creating
asymmetric losses that cascade through the correspondent banking system.
Under atomic settlement, both legs complete or neither does.

\textbf{Real-time supervisory visibility.} On-chain trade repositories
provide live aggregate exposure data, not T+1 or T+2 reporting. Position
concentration, limit utilisation, and collateral adequacy are visible in
real time, qualitatively superior to anything achievable under current
infrastructure.

\textbf{Immutable audit trail.} Every IM call, close-out trigger, and
collateral transfer is timestamped on-chain and cryptographically verifiable.
Regulatory model validation can use on-chain records rather than
firm-submitted data, permanently reducing the information asymmetry that
currently requires regulators to accept firms' characterisations of their
own models.

\textbf{Evidence base for MPoR reform.} The CFTC's August 2026 rulemaking
requires quantitative evidence that blockchain settlement genuinely compresses
close-out mechanics. On-chain data from executed bilateral derivatives
provides this evidence on a demonstrated rather than theoretical basis.

\textbf{Clearing mandate recalibration.} If blockchain bilateral derivatives
achieve operational equivalence with or superiority to cleared derivatives
in close-out speed, collateral certainty, and transparency, the policy
basis for mandatory clearing---the operational inferiority of bilateral
arrangements---is eroding. Regulators have an independent interest in
assessing whether the current clearing mandate remains optimally calibrated
for a world in which bilateral settlement can be atomic.

% ══════════════════════════════════════════════════════════════════════════════
\section{Regulatory Framework and Roadmap}
\label{sec:regulatory}
% ══════════════════════════════════════════════════════════════════════════════

The regulatory environment for blockchain-based derivatives collateral has
shifted materially since 2024:

\begin{itemize}[leftmargin=2em,itemsep=4pt]
  \item \textbf{CFTC Letter 25-39 (December 8, 2025):} Formal CFTC guidance
    that tokenised assets retain margin eligibility if ``functionally
    equivalent'' to their traditional form. BTC, ETH, and USDC pilot
    programme launched \citep{cftc_2025}.
  \item \textbf{GENIUS Act (July 2025):} Federal statutory framework for
    payment stablecoins, creating legally recognised instruments eligible for
    use as derivatives margin.
  \item \textbf{BCBS chair statement (November 2025):} Confirmed that growth
    of stablecoins and US/UK decisions not to implement the 1250\% public
    blockchain risk weight ``calls for a different approach.'' Formal BCBS
    review underway.
  \item \textbf{CFTC rulemaking target (August 2026):} Acting Chairman Pham
    outlined a timeline for technical amendments to CFTC regulations for
    margin, clearing, settlement, reporting, and recordkeeping to enable
    blockchain. This is the primary legislative vehicle for MPoR reduction
    recognition \citep{dwt_2025}.
  \item \textbf{ISDA position (2018):} ISDA stated that the ten-day bilateral
    MPoR ``should be reviewed'' once clearing volumes exceeded mandated
    levels \citep{isda_2018}. Those volumes have long been exceeded.
\end{itemize}

The MPoR reform argument is not merely policy preference. It follows
directly from the regulatory intent behind the ten-day floor: bilateral
derivatives were harder to close out than cleared ones. Section~\ref{subsec:m7}
demonstrates that this is no longer true for blockchain-confirmed bilateral
trades. If regulators accept the operational comparisons in that section---
and the CFTC August 2026 rulemaking provides exactly the forum to do so---the
appropriate regulatory outcome is not MPoR parity with the five-day cleared
floor but a shorter bilateral MPoR for fully on-chain portfolios. The EAD
reductions of equations~\eqref{eq:ead1}--\eqref{eq:ead01} show what one-day
and sub-day MPoR recognition would mean for regulatory capital.

% ══════════════════════════════════════════════════════════════════════════════
\section{Conclusion}
\label{sec:conclusion}
% ══════════════════════════════════════════════════════════════════════════════

This paper has provided a comprehensive quantitative analysis of how
blockchain trade confirmation and atomic settlement restructure the bilateral
derivatives cost stack across four dimensions. The primary analytical result
is that MPoR compression reduces SIMM risk weights, SA-CCR maturity factors,
and full-simulation PFE profiles by the same factor $\sqrt{N_2/N_1}$
simultaneously---establishing structural consistency across all three
frameworks through Propositions~1 and~2. For the regulatory five-day MPoR
recognition, this factor is $1/\!\sqrt{2} \approx 0.707$, producing a
29.3\% reduction that is invariant to portfolio composition, counterparty
credit quality, and netting set size.

Combined with seven secondary mechanisms, aggregate effects are: SIMM IM
down 40\%; all-in XVA down 60--66\%; SA-CCR capital down 68--84\%; credit
limit capacity up 3.6--8.6 times. These are not marginal refinements. They
are a structural repricing of bilateral derivatives risk available to any
institution that confirms trades on-chain and settles atomically.

The deeper argument is that the regulatory logic underlying the bilateral
MPoR penalty inverts under blockchain. Blockchain-confirmed bilateral
derivatives exhibit faster close-out than CCP default management procedures,
greater collateral certainty than CCP default waterfall exposure, zero
settlement risk compared to T+2 clearing settlement, and no systemic
concentration compared to too-big-to-fail central counterparties. If
regulators accept these operational comparisons---and the CFTC August 2026
rulemaking provides exactly the forum---the correct policy outcome is not
MPoR parity with the five-day cleared floor. It is a shorter bilateral
MPoR, potentially one day or less for fully on-chain portfolios. Equations
\eqref{eq:ead1} and \eqref{eq:ead01} show what that means for regulatory
capital. The question is not whether this will be recognised. It is when.

When Wall Street's 1968 paperwork crisis was resolved, the Depository Trust
Company did not improve the securities. It replaced the infrastructure
beneath them. What DTC achieved for equity settlement---dematerialisation
of paper into book entries, confirmation reduced to a computational event,
settlement made certain---blockchain achieves for bilateral derivatives.
The regulatory and pricing consequences are now quantifiable. The
infrastructure to realise them is operational. The work is to make the
argument clearly enough that regulators recognise what has already changed.

% ── BIBLIOGRAPHY ──────────────────────────────────────────────────────────────
\bibliographystyle{apalike}

\begin{thebibliography}{99}

\bibitem[BCBS(2014)]{bcbs_2014}
Basel Committee on Banking Supervision (2014).
\newblock \emph{The Standardised Approach for Measuring Counterparty Credit
  Risk Exposures} (CRE52).
\newblock Bank for International Settlements.

\bibitem[BCBS-IOSCO(2015)]{bcbs_iosco_2015}
Basel Committee on Banking Supervision and International Organization of
  Securities Commissions (2015).
\newblock \emph{Margin Requirements for Non-Centrally Cleared Derivatives}.
\newblock Final Rule. Bank for International Settlements / IOSCO.

\bibitem[Bank of England(2022)]{pra_2022}
Bank of England / Prudential Regulation Authority (2022).
\newblock PRA's review of the use of the SIMM model: Conclusions.
\newblock Letter from D.~Mackinnon, D.~Bailey, and N.~Benjamin, June 2022.

\bibitem[Burgard and Kjaer(2013)]{burgard_kjaer_2013}
Burgard, C. and Kjaer, M. (2013).
\newblock Funding strategies, funding costs.
\newblock \emph{Risk Magazine}, December 2013, 82--87.

\bibitem[CFTC(2025)]{cftc_2025}
Commodity Futures Trading Commission (2025).
\newblock \emph{Letter No.\ 25-39: Tokenized Collateral Guidance}.
\newblock December 8, 2025.

\bibitem[Davis Wright Tremaine(2025)]{dwt_2025}
Davis Wright Tremaine LLP (2025).
\newblock CFTC Issues Guidance on Tokenized Collateral and Other Crypto
  Sprint Developments.
\newblock December 2025.

\bibitem[Green and Kenyon(2014)]{green_kenyon_2014}
Green, A. and Kenyon, C. (2014).
\newblock KVA: Capital valuation adjustment.
\newblock \emph{Risk Magazine}, December 2014.

\bibitem[Gregory(2020)]{gregory_2020}
Gregory, J. (2020).
\newblock \emph{The XVA Challenge: Counterparty Risk, Funding, Collateral,
  Capital and Initial Margin}. 4th~ed.
\newblock Wiley Finance.

\bibitem[IOSCO(2023)]{iosco_2023}
International Organization of Securities Commissions (2023).
\newblock \emph{Streamlining Variation Margin Processes and Initial Margin
  Responsiveness of Margin Models in Non-Centrally Cleared Markets}.
\newblock Final Report. IOSCOPD787.

\bibitem[ISDA(2018)]{isda_2018}
International Swaps and Derivatives Association (2018).
\newblock Margin Rules and ISDA SIMM.
\newblock Position Paper, May 25, 2018.

\bibitem[ISDA(2022)]{isda_gov_2022}
International Swaps and Derivatives Association (2022).
\newblock \emph{ISDA SIMM Governance Framework}.
\newblock Version July 2022.

\bibitem[ISDA(2025a)]{isda_2025a}
International Swaps and Derivatives Association (2025a).
\newblock \emph{ISDA SIMM Methodology, Version 2.8+2506}.
\newblock Effective December 6, 2025.

\bibitem[ISDA(2025b)]{isda_2025b}
International Swaps and Derivatives Association (2025b).
\newblock \emph{ISDA SIMM: The Trusted Standard for Initial Margin
  Calculations}.
\newblock White Paper, November 2025.

\bibitem[Pykhtin and Zhu(2007)]{pykhtin_zhu_2007}
Pykhtin, M. and Zhu, S. (2007).
\newblock A guide to modelling counterparty credit risk.
\newblock \emph{GARP Risk Review}, July/August 2007, 16--22.

\bibitem[SIFMA(2025)]{sifma_2025}
Securities Industry and Financial Markets Association (2025).
\newblock Request for Input on the Use of Tokenized Collateral Including
  Stablecoins in Derivatives Markets.
\newblock Letter to CFTC, November 25, 2025.

\end{thebibliography}

\end{document}
